% =========================================================
% frontmatter/objective.tex
% =========================================================

\chapter*{Project Objective}
\thispagestyle{fancy}
\addcontentsline{toc}{chapter}{Project Objective}
\begingroup
\setlength{\emergencystretch}{2em}

The objective of this project is to develop and evaluate a complete
antenna-engineering workflow for a 5G microstrip patch antenna and a
two-element phased patch array operating near 2.501~GHz.

The work integrates closed-form antenna design, parameterized
full-wave electromagnetic simulation, impedance matching, radiation
analysis, mutual-coupling characterization, and phase-controlled beam
steering.

The project has four principal technical objectives:

\begin{enumerate}
    \item Design an inset-fed rectangular microstrip patch antenna using
    analytical transmission-line relations for the patch dimensions,
    effective dielectric constant, fringing-field extension,
    microstrip-feed width, and initial inset position.

    \item Implement and tune the single-patch model in CST Studio Suite,
    and evaluate its resonant frequency, reflection coefficient,
    directivity, radiation pattern, radiation efficiency, and total
    efficiency.

    \item Extend the tuned element into a two-port $2\times1$ patch array,
    quantify the effect of mutual coupling, and determine the separate
    tuning required to maintain acceptable input matching near the
    operating frequency.

    \item Demonstrate beam steering by applying controlled relative phase
    shifts between the two array elements and comparing the simulated
    scan angles with ideal two-element array-factor predictions.
\end{enumerate}

A secondary objective is to evaluate the numerical reliability and
physical limitations of the simulation model. Particular attention is
given to the CST Learning Edition mesh-cell restriction, the use of PEC
metallization in the baseline models, the supplementary
finite-conductivity copper sensitivity study, and the absence of formal
mesh- and boundary-convergence studies.

The final outcome is intended to provide a technically traceable
portfolio record that connects analytical antenna theory, full-wave
simulation, numerical interpretation, and phased-array behavior.
All reported performance results are analytical or simulated; no
fabricated prototype or measured antenna data are claimed.
\endgroup
